%! AP Statistics — Unit 2, Lesson 2.3 — Homework (Answer Key)
\documentclass[10pt]{article}
\usepackage{apstats-article}
\usepackage{apstats-key}

%=============================================================================
\begin{document}

\pageheader{Unit 2, Lesson 2.3}{Homework --- Answer Key}
\begin{tabularx}{\linewidth}{X X X}
Name:\;\ans{Key} & Date:\; & Period:\;
\end{tabularx}

\vspace{0.2in}

% ── PROBLEM 1 ─────────────────────────────────────────────────────────────────
\begin{tcolorbox}[colback=sky, colframe=navy, boxrule=0.9pt, arc=2mm,
    left=4mm, right=4mm, top=3mm, bottom=3mm]
\small\textbf{Problem 1 --- Helmet Use and Injuries}

\vspace{8pt}
\renewcommand{\arraystretch}{1.8}
\begin{tabularx}{\linewidth}{@{} l C{3cm} C{3cm} C{2.4cm} @{}}
\rowcolor{sky}
                          & \textbf{Injury: Yes} & \textbf{Injury: No} & \textbf{Row Total} \\
\hline
Helmet: Yes      & 15 & 285 & \ans{300} \\
Helmet: No       & 40 & 160 & \ans{200} \\
\hline
Col.\ Total      & \ans{55} & \ans{445} & \ans{500} \\
\end{tabularx}
\end{tcolorbox}

\vspace{0.14in}

\begin{enumerate}[leftmargin=1.6em, itemsep=0.16in]

    \item $\hat{p}_{\text{helmet}} = \ans{15/300 = 0.05}$
          \hfill $\hat{p}_{\text{no helmet}} = \ans{40/200 = 0.20}$

    \item Difference: \ans{$0.05 - 0.20 = -0.15$}\\[8pt]
          \ansline{Helmet wearers had a head-injury rate 15 percentage points \emph{lower}
          than non-wearers (5\% vs.\ 20\%).}

    \item RR: \ans{$0.05/0.20 = 0.25$}\\[8pt]
          \ansline{Helmet wearers were 0.25 times as likely to sustain a head injury as
          non-wearers --- i.e., helmet wearers had only one-quarter the injury risk of non-wearers.}

    \item \ansline{The headline is not accurate. The actual relative risk is 0.25 (helmet vs.\
          no helmet), meaning non-wearers are $1/0.25 = 4$ times as likely to be injured ---
          \emph{four} times the rate, not twice. The headline underestimates the association.}

\end{enumerate}

\vspace{0.22in}

% ── PROBLEM 2 ─────────────────────────────────────────────────────────────────
\begin{tcolorbox}[colback=sky, colframe=navy, boxrule=0.9pt, arc=2mm,
    left=4mm, right=4mm, top=3mm, bottom=3mm]
\small\textbf{Problem 2 --- Screen Time and Sleep Quality}
\end{tcolorbox}

\vspace{0.14in}

\begin{enumerate}[leftmargin=1.6em, itemsep=0.16in]

    \item
    \renewcommand{\arraystretch}{1.8}
    \begin{tabularx}{\linewidth}{@{} l C{2.8cm} C{2.8cm} C{2.4cm} @{}}
    \rowcolor{sky}
                              & \textbf{Sleep: Poor} & \textbf{Sleep: Good} & \textbf{Row Total} \\
    \hline
    Screen: High     & \ans{63} & \ans{27} & \ans{90} \\
    Screen: Low      & \ans{45} & \ans{105} & \ans{150} \\
    \hline
    Col.\ Total      & \ans{108} & \ans{132} & \ans{240} \\
    \end{tabularx}

    \item $\hat{p}_{\text{high}} = \ans{63/90 = 0.70}$
          \quad $\hat{p}_{\text{low}} = \ans{45/150 = 0.30}$\\[10pt]
          Difference: \ans{$0.70 - 0.30 = 0.40$}
          \hfill RR: \ans{$0.70/0.30 \approx 2.33$}

    \item \ansline{Among teenagers with high screen time before bed, 70\% reported poor sleep,
          compared to only 30\% of those with low screen time --- a difference of 40 percentage
          points. High-screen-time teens were about 2.33 times as likely to report poor sleep,
          suggesting a strong positive association between screen time and poor sleep quality.}

    \item \ansline{For the helmet problem, the outcome (head injury, 5--20\%) is relatively
          rare, so the relative risk (4$\times$) is more striking than the 15-pp absolute
          difference. For screen time, the outcome is common (30--70\%), so the absolute
          difference (40 pp) is also very interpretable. The relative risk is especially
          useful when outcomes are rare, because a small absolute difference can still
          represent a dramatic multiplicative change.}

\end{enumerate}

\begin{teachernote}
\textbf{Problem 1, item 4:} Many students confuse ``twice as likely'' with a RR of 2.0.
Here the RR (non-helmet / helmet) $= 0.20/0.05 = 4.0$. The headline is wrong by a factor of 2.
Use as a media-literacy discussion: statistics can be selectively or inaccurately reported.

\textbf{Problem 2, reflection:} Accept any well-reasoned answer. The key idea is that the
appropriateness of each statistic depends on the prevalence of the outcome.
\end{teachernote}

\end{document}
